%%%%%%%%%%  scan idx 54  |  folio 41  %%%%%%%%%%
\SGAPrintPage{41}%
\sgaStableTarget{sga6.e00.eq.2-8}\begin{equation}\tag{2.8}
 C^1(L(D))=\ellc(D)\qquad C^i(L(D))=0\quad (i>1)
\end{equation}
(where,
for an arbitrary cycle $Z$, we denote by $\ellc(Z)$ its class in $A(X)$). Finally, we require the $C^i$ to be \emph{functorial}, i.e. that, when $f$ is a morphism from $X$ to $Y$ and $E$ is a vector bundle on $Y$, we have
\sgaStableTarget{sga6.e00.eq.2-9}\begin{equation}\tag{2.9}
 C^i(f^!(E))=f^*(C^i(E)),
\end{equation}
where $f^!(E)$ is the inverse image of $E$ on $X$. We can also write (2.9) in the following form: the notion of inverse image of a bundle defines a $\lambda$-homomorphism of augmented $\lambda$-rings
\sgaStableTarget{sga6.e00.eq.2-10}\begin{equation}\tag{2.10}
 f^!:\underline{K}(Y)\longrightarrow \underline{K}(X)\SGAAuthorFootnote{*}{The notation $f^!$ used here conflicts with that used in duality theory (cf.\ for example R. Hartshorne, Residues and Duality, Lecture Notes in Math. no.\ 20, Springer, 1966). This is why in this Seminar we replace it by the notation $f^*$.}
\end{equation}
(here we write $\underline{K}(X)$ instead of $\underline{K}(\xi(X))$), and (2.9) means that the diagram commutes:
\[
\begin{tikzcd}[column sep=large,row sep=large]
K(Y) \arrow[r,"f^!"] \arrow[d,"\Ctil_Y"'] & K(X) \arrow[d,"\Ctil_X"]\\
\Atil(Y) \arrow[r,"f^*"] & \Atil(X).
\end{tikzcd}
\]
(The space on which the symbol $\Ctil$ is considered has been put as a subscript.)

\sgaStableTarget{sga6.e00.thm.2-2}\textbf{Theorem 2.2.} \emph{There exists a theory of Chern classes satisfying conditions (2.5), (2.7), (2.8) and (2.9).}

We shall also admit the following proposition:

\sgaStableTarget{sga6.e00.prop.2-3}\textbf{Proposition 2.3.} \emph{Let $Y$ be a closed subset of $X$ and $U$ its complement; if $x$ is an element of $A(X)$ inducing $0$ on $U$, there exists a representative of $x$ that is a cycle supported on $Y$.}

%%%%%%%%%%  scan idx 55  |  folio 42  %%%%%%%%%%
\SGAPrintPage{42}%
\sgaStableTarget{sga6.e00.cor.unnamed.occ2}\textbf{Corollary.} \emph{Let $p=\dim X-\dim Y$. Then, if $i<p$, the homomorphism $A^i(X)\longrightarrow A^i(U)$ is bijective. Every element of the kernel of $A^p(X)\longrightarrow A^p(U)$ is a linear combination of components of dimension $n-p$ of $Y$ (and is therefore proportional to $\ellc(Y)$ if $Y$ is irreducible)}

\sgaStableTarget{sga6.e00.sec.2.occ3}\subsection*{\S 2. Definition of the Chern classes of coherent algebraic sheaves}

Let $\mathbb{C}$ be an abelian category, and let $\mathbb{C}_0$ be a subclass of $\mathbb{C}$. We denote by $K(\mathbb{C}_0)$ the quotient group of the free group generated by the classes (up to isomorphism) of elements of $\mathbb{C}_0$, by the subgroup generated by elements of the form $(E)-(E')-(E'')$, where $E$ is an \emph{extension} of $E'$ by $E''$ ($E,E',E''\in\mathbb{C}_0$). We must assume that the classes of elements of $\mathbb{C}_0$ for the isomorphism relation form a set. We denote by $\gamma_{\mathbb{C}_0}(E)$ the class of an $E\in\mathbb{C}_0$ in $K(\mathbb{C}_0)$, and omit the subscript $\mathbb{C}_0$ when there is no risk of confusion. The group $\underline{K}(\mathbb{C}_0)$ is the solution of a \emph{universal problem} concerning functions $E\longrightarrow f(E)$ defined on $\mathbb{C}_0$, with values in an arbitrary abelian group $A$, which are \emph{additive} in the sense that $f(E)=f(E')+f(E'')$ whenever $E$ is an extension of $E'$ by $E''$.

Let $F$ be an additive functor from $\mathbb{C}$ to $\mathbb{C}'$, and $\mathbb{C}'_0$ a subclass of $\mathbb{C}'$ satisfying the same condition as $\mathbb{C}_0$, such that, for every $A\in\mathbb{C}_0$, the $R^qF(A)$ are in $\mathbb{C}'_0$ and are zero for sufficiently large $q$ (we assume that injective resolutions exist in $\mathbb{C}$, so that the derived functors $R^qF$ of $F$ exist); then $F$ defines a homomorphism from $\underline{K}(\mathbb{C}_0)$ to $\underline{K}(\mathbb{C}'_0)$, still denoted $F$, by the formula
\sgaStableTarget{sga6.e00.eq.2-11}\begin{equation}\tag{2.11}
 F(\gamma_{\mathbb{C}_0}(E))=\sum_{i\geq 0}(-1)^i\gamma_{\mathbb{C}'_0}(R^iF(E)).
\end{equation}
More generally, when we have a cohomological functor $(F^i)$ from $\mathbb{C}$ to $\mathbb{C}'$
%%%%%%%%%%  scan idx 56  |  folio 43  %%%%%%%%%%
\SGAPrintPage{43}%
such that,
for every $E\in\mathbb{C}_0$, the $F^i(E)$ are in $\mathbb{C}'_0$ and are zero except for finitely many, we deduce a homomorphism from $\underline{K}(\mathbb{C}_0)$ to $\underline{K}(\mathbb{C}'_0)$. These remarks extend to multifunctors; moreover, the ``classical'' spectral sequences give transitivity properties that we shall not make explicit.

In particular, when $F$ is an exact functor, formula (2.11) simplifies:
\sgaStableTarget{sga6.e00.eq.2-11bis}\begin{equation}\tag{2.11 bis}
 F(\gamma_{\mathbb{C}_0}(E))=\gamma_{\mathbb{C}'_0}(F(E)).
\end{equation}
This applies in particular when $\mathbb{C}=\mathbb{C}'$, when $F$ is the identity functor, and when $\mathbb{C}_0\subset\mathbb{C}'_0$: we have a canonical homomorphism from $\underline{K}(\mathbb{C}_0)$ to $\underline{K}(\mathbb{C}'_0)$.

\sgaStableTarget{sga6.e00.thm.2-3}\textbf{Theorem 2.3.} \emph{Let $\mathbb{C}$ be an abelian category and $\mathbb{C}_0\subset\mathbb{C}_1$ two subclasses such that the classes} (up to isomorphism) \emph{of objects of $\mathbb{C}_1$ form a set. We assume that the following conditions hold}
\begin{enumerate}[label=(\roman*)]
\item \emph{For every exact sequence $0\to A'\to A\to A''\to0$ in $\mathbb{C}$ such that $A$ and $A''$ are in $\mathbb{C}_0$ (resp.\ $\mathbb{C}_1$), we have $A'\in\mathbb{C}_0$ (resp.\ $A'\in\mathbb{C}_1$).}
\item \emph{Let $u:A\to B$ and $v:C\to B$ with $A,B,C\in\mathbb{C}_1$ and $u$ surjective; then there exist $L\in\mathbb{C}_0$ and homomorphisms $u':L\to C$ and $v':L\to A$ such that $u'$ is surjective and $vu'=uv'$.}
\item \emph{For every $A\in\mathbb{C}_1$, there exists an integer $d=d(A)$ such that, for every left resolution $\underline{L}$ of $A$ by objects of $\mathbb{C}_0$, we have $Z_d(\underline{L})\in\mathbb{C}_0$} ($Z_d$ denotes the cycles of dimension $d$).
\end{enumerate}
\emph{Under these conditions, the canonical homomorphism from $K(\mathbb{C}_0)$ to $K(\mathbb{C}_1)$ is an isomorphism.}

\textbf{Principle of the proof:} let $A\in\mathbb{C}_1$; by (ii) and (iii) there exists a finite resolution $\underline{L}$ of $A$ by objects of $\mathbb{C}_0$; put
%%%%%%%%%%  scan idx 57  |  folio 44  %%%%%%%%%%
\SGAPrintPage{44}%
then $f(A)=\sum_{i\geq 0}(-1)^i\gamma_{\mathbb{C}_0}(L^i)$.
We show that the element of $\underline{K}(\mathbb{C}_0)$ thus constructed does not depend on the chosen resolution $\underline{L}$: given a second resolution $\underline{L}'$, we ``cover'' $\underline{L}$ and $\underline{L}'$ by a third resolution $\underline{L}''$, with \emph{surjective} homomorphisms from $\underline{L}''$ to $\underline{L}$ and $\underline{L}'$ (which is possible by (i) and (ii)); we are then reduced to proving $\underline{E}_{\mathbb{C}_0}(\underline{L})=\underline{E}_{\mathbb{C}_0}(\underline{L}'')$. Now the difference of the two is $\underline{E}_{\mathbb{C}_0}(\underline{N})$, where $\underline{N}$ is the kernel of the surjective homomorphism from $\underline{L}''$ onto $\underline{L}$, a kernel that is in $\mathbb{C}_0$ by (i). Now $\underline{N}$ is an acyclic complex in all dimensions, whence $\underline{E}_{\mathbb{C}_0}(\underline{N})=0$, whence the conclusion (N.B. $\underline{E}_{\mathbb{C}_0}(\underline{L})$, for a finite complex with values in $\mathbb{C}_0$, denotes the Euler--Poincaré characteristic $\sum_{i\geq0}(-1)^i\gamma_{\mathbb{C}_0}(L^i)$ of $\underline{L}$).

We prove analogously that the function $f:\mathbb{C}_1\longrightarrow\underline{K}(\mathbb{C}_0)$ thus constructed is \emph{additive}, and therefore defines a homomorphism from $\underline{K}(\mathbb{C}_1)$ to $\underline{K}(\mathbb{C}_0)$, and it is immediate that this homomorphism is inverse to the canonical homomorphism from $\underline{K}(\mathbb{C}_0)$ to $\underline{K}(\mathbb{C}_1)$. q.e.d.

\sgaStableTarget{sga6.e00.rem.unnamed.occ3}\textbf{Remark:} To verify (ii), it suffices to verify that $\mathbb{C}_1$ is contained in a subclass $\mathbb{C}_2$ of $\mathbb{C}$ satisfying the two conditions
\begin{enumerate}[label=(ii \alph*)]
\item Every $A\in\mathbb{C}_2$ is isomorphic to a quotient of an $L\in\mathbb{C}_0$.
\item If $u:A\to B$, with $A,B\in\mathbb{C}_2$, then $\Ker u\in\mathbb{C}_2$.
\end{enumerate}

\sgaStableTarget{sga6.e00.cor.unnamed.occ3}\textbf{Corollary:}
\emph{Let $X$ be a quasi-projective algebraic space. Let $\cF(X)$ be the category of coherent algebraic sheaves on $X$, $\xi(X)$ the category of locally free coherent algebraic sheaves} (equivalent to the category of vector bundles on $X$), \emph{and $\cF_0(X)$ the category of coherent algebraic sheaves $F$ such that, for every $x\in X$, $F_x$ is an $\underline{O}_x$-module
%%%%%%%%%%  scan idx 58  |  folio 45  %%%%%%%%%%
\SGAPrintPage{45}%
of finite cohomological dimension (we therefore have $\xi(X)\subset\cF_0(X)\subset\cF(X)$). Then the canonical homomorphism}
\sgaStableTarget{sga6.e00.eq.2-12}\begin{equation}\tag{2.12}
 K(\xi(X))\longrightarrow K(\cF_0(X))
\end{equation}
\emph{is an isomorphism.}

We are going to verify the conditions of theorem 2.3. Conditions (i) and (iii) follow from the definition of cohomological dimension and from the fact that a projective module of finite type over the noetherian local ring $\underline{O}_x$ is free. To verify (ii), we shall show that $\mathbb{C}=\cF(X)$ satisfies conditions (ii a) and (ii b) of the remark. The second is trivial; it therefore remains to prove the following result:

\sgaStableTarget{sga6.e00.prop.2-4}\textbf{Proposition 2.4. (Serre).} \emph{Let $X$ be a quasi-projective algebraic space. Then every coherent algebraic sheaf on $X$ is isomorphic to the quotient of a locally free algebraic sheaf.}

This proposition holds when $X$ is a closed subset of projective space, by [2], since, for sufficiently large $n$, $F(n)$ is generated by its sections, which means that $F$ is a quotient of $\underline{O}(-n)^k$. In the general case, $X$ is an open subset of a closed subset of a projective space, and it suffices to use the following result (due to Cartier and Serre in particular cases):

\sgaStableTarget{sga6.e00.prop.2-5}\textbf{Proposition 2.5.} \emph{Let $F$ be a coherent algebraic sheaf on an open subset $U$ of an algebraic space $X$. Then $F$ is isomorphic to the restriction of a coherent algebraic sheaf on $X$.}

By applying Zorn's lemma to the open subsets containing $U$ to which $F$ can be extended, we are reduced to the case where $X$ is affine. Let $\overline F$ be the
alge%
